\begin{otsblock}[Gewisslich sterben]
        Wie schon geschrieben, müssen wir alle eines Tages sterben. Da kommen wir nicht drumherum. Jeder, ob arm, reich, nett oder böse geht den gleichen Weg.  Das ist der Beweis, dass das Gebot von Gott nicht eingehalten wurde. Wie das passiert ist, könnt ihr in der Bibel in \biblezit{IMos}{3:1-6} nachlesen. Aber was ist denn passiert?
        \begin{bibelbox}{ELB}{IMos}{3:7}
            \hh{7}Da wurden ihre beiden Augen aufgetan, und sie erkannten, dass sie nackt waren; und sie hefteten Feigenblätter zusammen und machten sich Schurze.
        \end{bibelbox}
             
        \begin{bibelbox}{ELB}{IMos}{1:31}
            \hh{31} Und Gott sah alles, was er gemacht hatte, und siehe, es war sehr gut. Es wurde Abend, und es wurde Morgen: der sechste Tag
        \end{bibelbox}
        Gott fand also alles, was er geschaffen hatte als sehr gut. Gott weiss genau, was die eben erschaffenen Menschen im Schilde führen. Wenn wir das so lesen, denken wir oft: \enquote{Bei mir wäre das nicht passiert!} Wirklich?

        Wir sind ja hier alle schon ein bisschen älter. Kennt ihr diesen Trickfilm La Linea? Wieso hat Gott es mit dem Menschen nicht so gemacht wie der Zeichner bei diesem Trickfilm? Einfach drüberwischen und alles wieder löschen.
        \begin{itemize}
            \item Gott erschafft Himmel und Erde mit dem Menschen
            \item Der Mensch sündigt
            \item Gott hat einen Plan
        \end{itemize}
        
        Diesen Plan Gottes lesen wir in 1Mos 3:15. 
        \begin{bibelbox}{SCHL}{1Mos}{3:15}
            Und ich werde Feindschaft setzen zwischen dir und der Frau, zwischen deinem Nachwuchs und ihrem Nachwuchs; er wird dir den Kopf zermalmen, und du, du wirst ihm die Ferse zermalmen.
        \end{bibelbox}
        In diesem Vers wird der ganze Plan Gottes aufgezeigt. Er gibt uns nicht auf. Er hat uns geschaffen um mit uns Gemeinschaft zu haben. Wenn Gott etwas macht, dann bringt er es auch zu Ende. Aber, Gott macht es so wie er will und nicht wie wir uns das Denken.

        Über diesen Vers sind ganze Bücher geschrieben worden. Eva wusste was Gott mit diesen Worten meinte, denn sie war dauernd in Erwartung auf diesen Erlöser, der, welcher der Schlange den Kopf zermalmt. 
    \end{otsblock} 
    \begin{otsblock}[ABER]
        Der Mensch musst das Paradies verlassen und unter Mühsal und Schweiss sein Unterhalt verdienen. Die Frau kann nur unter schmerzen gebären. Ich weiss nicht genau, aber ich habe viele Kühe kalbern sehen. Bei denen wo es keine komplikationen gab, sah das nicht aus als ob die Kühe schmerzen hatten. Sie frassen sogar wären dem Pressen noch Heu. Wenn das Kalb am Boden lagen, drehten sie sich kauend zu ihm hin. Oder Samuel? Hast du das auch so erlebt?
        
        Gott hat ihr Umfeld verändert, aber nicht den Menschen. Der Mensch darf immer noch frei und selbst bestimmen und entscheiden. \betonung{Gott will Menschen, die Freiwillig zu ihm kommen.} Aber die ganze Schöpfung ist dem Untergang geweiht. Es geht unaufhörlich Berg ab mit dieser Schöpfung. Nichts hat Bestand. Auch wenn uns die Evolutionstheoretiker erklären wollen, dass sich alles weiter entwickelt, sehen wir doch allzu gut, dass auf dieser Welt das Gegenteil passiert. Letzthin habe ich ein Foto gesehen, wo erwachsenen Männer mit Schnuller und Windeln in einem Umzug mit liefen. Die Krönung der Schöpfung.

       So lesen wir in:
        \begin{bibelbox}{ELB}{Rom}{8:20}
            Denn die Schöpfung ist der Nichtigkeit unterworfen worden - nicht freiwillg, sondern durch den, der sie unterworfen hat -- auf Hoffnung hin.
        \end{bibelbox}
        Oder:
        \begin{bibelbox}{NUE}{Rom}{8:20}
            Denn alles Geschaffene ist der Vergänglichkeit ausgeliefert -- undfreiwillig. Gott hat es so verfügt.
        \end{bibelbox}
        Und wie das aussieht lesen wir in:
        \begin{bibelbox}{ELB}{Rom}{8:22}
            Denn wir \betonung{wissen}, dass die \betonung{ganze} Schöpfung zusammen seufzt und zusammen in Geburtswehen liegt bis jetzt.
        \end{bibelbox}
    \end{otsblock}
    \begin{otsblock}
        Gott hat uns Menschen nicht aufgeben. Nein. Er führt seinen Plan mit uns zum Ende. Heute wissen wir wie er endet. Damals wussten Adam und Eva das noch. Sie lebten knapp tausend Jahre. Wie oft haben sie sich gegenseitig wohl Vorwürfe gemacht?
    \end{otsblock}
    \begin{otsblock}
        Nächstes Mal schauen zu sammen wie es weiter geht. Wir kommen dann zu Phase 3. Diese Phase ist eine lange Zeit und dauert bis zum ersten kommen von Jesus Christus unserem Retter. Und was in dieser Zeit Gott alles mit den Menschen angestellt, hat wird uns dann wohl den Rest des Jahres beschäftigen.
    \end{otsblock}    